\renewcommand{\refname}{7. Tài liệu tham khảo}
\addcontentsline{toc}{section}{7. Tài liệu tham khảo}

\begin{thebibliography}{9}

\bibitem[Anderson \& Krathwohl(2001)]{anderson2001}
Anderson, L. W., \& Krathwohl, D. R. (Eds.). (2001). \textit{A taxonomy for learning, teaching, and assessing: A revision of Bloom's taxonomy of educational objectives}. Longman.

\bibitem[Bates(2019)]{bates2019}
Bates, A. W. (2019). \textit{Teaching in a digital age: Guidelines for designing teaching and learning}. Tony Bates Associates Ltd.

\bibitem[Jonassen(1991)]{jonassen1991}
Jonassen, D. H. (1991). Evaluating constructivistic learning. \textit{Educational Technology}, 31(9), 28--33.

\bibitem[Mishra \& Koehler(2006)]{mishra2006}
Mishra, P., \& Koehler, M. J. (2006). Technological pedagogical content knowledge: A framework for teacher knowledge. \textit{Teachers College Record}, 108(6), 1017--1054.

\bibitem[Mollick \& Mollick(2023)]{mollick2023}
Mollick, E. R., \& Mollick, L. (2023). \textit{Using AI to implement effective teaching strategies in classrooms: Five strategies for educators} (Wharton School Research Paper). SSRN. \url{https://doi.org/10.2139/ssrn.4391243}

\bibitem[Phan(2023)]{phan2023}
Phan, T. T. T. (2023). Ứng dụng trí tuệ nhân tạo (AI) trong giáo dục đại học tại Việt Nam: Cơ hội và thách thức. \textit{Tạp chí Khoa học Giáo dục Việt Nam}, 19(2), 15--20.

\bibitem[Siemens(2005)]{siemens2005}
Siemens, G. (2005). Connectivism: A learning theory for the digital age. \textit{International Journal of Instructional Technology and Distance Learning}, 2(1), 3--10.

\end{thebibliography}
